% Integrated 2026-09-28 from pipeline/drafts/2026-09-28/related-work-v4.tex.
% Source: research/literature/positioning-a-star-r/dossier.md (verdict: partial
% overlap). Markers R23--R63 map to rows in claim-evidence.md, section
% "Related work v4 (integrated 2026-09-28)". R01--R22 describe the superseded
% section and are retired. No v4 gate or tier result is stated or implied.

% R23
Gestefeld and Lorenz calibrate an interaction model to European Social Survey distributions and to 50,008 annual opinion transitions from 10,627 SHP respondents \cite{gestefeld2023}.
% R24
Their macro-level fit to opinion distributions uses a binomial reference, whereas their micro-level evaluation compares empirical and simulated transition-probability matrices by Euclidean distance \cite{gestefeld2023}. % R24 RC51
Neither comparison is the matched persistence null used here. They also note that the absence of individual communication records puts some questions beyond the reach of these data \cite{gestefeld2023}. % R24 RC52
% R25
Using repeated individual surveys and observed communication events, the CoDiNG study finds that a baseline in which everyone keeps their opinion reaches a high overall F1 while detecting no change. It concludes that opinion-dynamics models often fail to outperform trivial baselines on real-world data \cite{coding2024}.
% R26
Achitouv and colleagues contrast interaction and persistence mechanisms using observed retweet exposures \cite{achitouv2026}.
% R27
Moor-Smith and Carpentras compare four opinion-dynamics models with a ``no change'' baseline on synthetic data and European Social Survey series. They report that a minimum amount of opinion drift is required before a model's explained variance exceeds that baseline \cite{moor2026}.
% R28
Adding noise to locate where reconstruction becomes infeasible, they find signs that the optimizer learns to mimic the null and ``freeze'' where the signal becomes unclear \cite{moor2026}.
% R29
Among the works we read, this is the closest antecedent: it asks how much drift a model needs to beat a no-change baseline on repeated cross-sections, whereas $a^*(r)$ asks how much drift individual panel trajectories need before they separate from matched persistence.
% R30
They use cross-sectional distributions as a proxy for longitudinal ground truth and add noise to the evolving process \cite{moor2026}; we analyze individual trajectories, place error only in a response channel whose signal share $r$ is fixed externally, match persistence nulls on observed lag-one (and, for two-parameter nulls, lag-two) correlations, and evaluate panel-level AUC with a fixed discriminator.

% R31
Carpentras and Quayle find that random noise and scale distortions, unlike binning, can strongly affect the predictions of the Deffuant model \cite{carpentras2022}. They also find that monotone scale distortions can make data generated by that model look incompatible with it, so poor predictions are not necessarily evidence of a poor model \cite{carpentras2023}.
% R32
In the parts of these studies we read, neither compares an influence model with a no-influence null or varies reliability in a panel.
% R33
Survey research frames the same ambiguity as real change versus response instability: Hill and Kriesi's four-wave Swiss panel analysis separates stable, unstable and changing respondents and reports results that undermine the measurement-error model \cite{hill2001,hillkriesi2001}.
% R34
Kiley and Vaisey note that if active updating is frequent enough for people to change between every wave, it becomes indistinguishable from settled dispositions with high measurement error \cite{kiley2020}.
% R35
Lersch shows with simulated data that persistent change cannot be reliably identified in a three-wave panel, and turns to long household panels, the SHP among them \cite{lersch2023}.
% R36
We therefore separate the response channel from latent updating and include trait-plus-autoregressive and mover--stayer structures among the persistence nulls.
% R37
The grid value $r$ is a declared sensitivity parameter rather than an estimate for the SHP trust item.
% R72
Political trust also has a measurement literature. Turper and Aarts \cite{turper2017} use latent autoregressive models to separate measurement error in political-trust items. Lugtig \cite{lugtig2017} uses a multi-trait multi-method model to estimate measurement error in political-trust questions and compare it with attrition error.
% R73
Those estimates concern different items, populations, scales and error models; they neither identify the simulator's $r$ nor supply a reliability value that can be transported to the SHP item.
% R64
Devine and Valgar{\dh}sson study stability and change in political trust across six panels, including the SHP federal-government trust item, using response changes, correlations and within--between models \cite{devine2024trust}.
% R65
They discuss single-item measurement error but do not estimate the SHP item's response reliability \cite{devine2024trust}.

% R38
Mulder presents simulation-based power analysis for the random-intercept cross-lagged panel model in the powRICLPM package, whose documentation defines conditions by sample size, number of waves, proportion of between-unit variance and indicator reliability \cite{mulder2023}.
% R39
With that tool, reliability of 0.60, 0.70 or 0.80, three to five waves and up to 10,000 respondents, Dang and Ma show that power to detect within-person cross-lagged effects depends jointly on reliability, intraclass correlation, number of waves, sample size and target effect size \cite{dang2026}.
% R40
Lucas shows in simulated data that, under realistic assumptions, the cross-lagged panel model is very likely to find spurious cross-lagged effects when they do not exist, which motivates our trait-plus-autoregressive null \cite{lucas2023}; Hamaker and colleagues also critique the same model \cite{hamaker2015critique}.
% R60
Trait--state--error models separate stable trait, changing state and error components in multiwave panel data \cite{kenny1995tse}.
% R61
Separating these components can be difficult. Cole and colleagues report improper solutions in simulations of trait--state--error models with few waves or very high or low occasion-factor stability and propose a trait--state--occasion model. L\"udtke and colleagues propose a Bayesian estimator to stabilize estimation of the STARTS model \cite{cole2005tso,ludtke2018starts}.
% R66
Schuurman and Hamaker separate measurement error from innovations in multilevel autoregressive models and show that ignoring measurement error can distort autoregressive and cross-lagged effects, including creating spurious paths \cite{schuurman2019reliability}.
% R67
Their examples concern dynamics among psychological variables within a person, rather than a test of interpersonal influence against matched persistence generators \cite{schuurman2019reliability}.
% R62
With the response channel added, our confirmatory null, trait-plus-autoregressive persistence (B), is a rank-transported trait--state--error process. It is given the true signal share $r$ and cannot adjust the response-channel noise variance, so the boundary measures detectability at a known signal share.
% R41
For influence in observed networks, Stadtfeld and colleagues show that power in stochastic actor-oriented models is strongly affected by network size, waves, effect sizes, missing data and turnover; with two waves and 30 actors, power for the influence effect is only 10 percent \cite{stadtfeld2020}.
% R42
Agent-based ground truth has also been used to test whether longitudinal estimators separate influence from homophily, without a response-error channel of the present kind \cite{zachrison2016,daza2021}.
% R43
Two further studies treat measurement error in social-influence inference as error in the observed network, not the response: Xu and Frank propose a simulation-based sensitivity analysis for selection mechanisms that bias network observations, and Advani and Malde review identification of social effects when networks form endogenously or are imperfectly observed \cite{xu2021sensitivity,advani2018credibly}.
% R44
Our study resembles an applied power analysis but differs in generator, null and criterion: it maps where a fixed discriminator separates two synthetic laws, rather than the power of a test applied to one observed panel.
% R45
The generator adds a nonlinear partner drift to an individual autoregression rather than a within-person linear cross-lagged path or an influence parameter on an observed network. The comparator is a family of no-network null generators matched on observed moments rather than a zero coefficient. The criterion is the AUC of a fixed discriminator on lag-three to lag-five features rather than the significance of an estimated parameter.
% R46
That more waves or a larger signal share ease detection is expected from this literature and is not a result we claim.
% R68
Soni, Ramirez and Eisenstein test influence by comparing the held-out accuracy of event-cascade rankers with and without influence features, and vary influence strength in simulations with self-excitation and homophily \cite{soni2019rankers}.
% R69
Their tests use a known network and event histories, so they do not provide a signal-share boundary for noisy ordinal panels \cite{soni2019rankers}.

% R47
Lenti, Monti and De Francisci Morales estimate bounded-confidence parameters by likelihood with fully observed opinions and interactions, partially observed interactions, and noisy opinion proxies \cite{lenti2024}; Borile and colleagues show that joint estimation is affected by identifiability issues in specific regions of parameter space \cite{borile2026}.
% R70
Lang and colleagues jointly infer an unknown weighted network and interaction kernel from multiple trajectories, and separately vary process forcing and observation noise in numerical experiments \cite{lang2026joint}.
% R71
Their target is network and kernel estimation error, rather than a minimum influence strength for discriminating panels from matched persistence nulls \cite{lang2026joint}.
% R48
Tang, Snijders and Flache separate bounded confidence from negative influence with a stochastic actor-oriented model on longitudinal friendship-network data \cite{tang2025empirical}.
% R49
Ureña-Carrion, Peixoto and Iñiguez assume a known network and fully observed node states, compare the detectability of binary spreading mechanisms, including an independent one, across parameter space, and show that statistical model selection can fail under common conditions \cite{urena2026}.
% R50
They present this, to their knowledge, as the first systematic effort to assess detectability limits and the prevalence of paradigmatic mechanisms in networked spreading \cite{urena2026}.
% R51
Our question shares the idea of a detectability map but not its observation regime: opinions are continuous and recorded through a noisy ordinal channel, the network is unobserved, and the discriminator is a fixed summary classifier rather than an exact likelihood.
% R52
Manski's reflection problem limits inference about endogenous peer effects without prior information on reference-group composition \cite{manski1993}, so any separation is read as discrimination between dynamic patterns.

% R53
Gelastopoulos, Sage and van de Rijt prove that any longitudinal data with systematic differences can also be produced by a reinforcement-driven process, so that pure reinforcement and pure heterogeneity models are generically confounded \cite{gelastopoulos2025}.
% R54
Their reinforcement is within-person, not interpersonal, but the result warns that trait-based nulls can mimic dynamics. The boundary $a^*(r)$ is a finite contrast against a fixed family matched at most at lags one and two, with the signal sought at lags three to five, so it does not contradict generic equivalence. Non-attainment is read as compatibility with the declared nulls, never as indistinguishability.
% R55
For each successful validated fit and training realization, the Bayes comparator knows the same $r$ and fitted null law as the fixed linear discriminant. Its optimal population AUC bounds that of the trained discriminant before averaging over the same successful-fit procedure. The resulting average of conditional optima is not the optimum against a marginal mixture of fitted nulls. Thus a population boundary attained on a complete valid grid upper-bounds the corresponding boundary of this conditional optimal comparator. The bound concerns population quantities rather than the estimated $\hat a^*(r)$ or its inverted interval; it is neither a minimax or asymptotic limit nor a test with unknown $r$ and null parameters, and it is probably loose, because an optimal procedure could use every person's trajectory rather than three pooled correlations.

% R56
We declared ten English-language queries on 27 September 2026 and ran them in one session on 27--28 September 2026. We used the queries without rewording on OpenAlex and Semantic Scholar; because the arXiv export API returned HTTP 406, we used their pre-declared Boolean translations through the arXiv site search. We screened at most the first 20 results per query and service, 399 titles in all. Eight records were relevant: five were records of four works discussed above and three were new \cite{carpentras2022,xu2021sensitivity,advani2018credibly}.
% R63
The per-run record was reconstructed from the session transcript; it reproduces these totals but lacks URLs, most run times and some identifiers (\hyperref[supp:search-record]{Supplementary Material: positioning search record}).
% R57
Coverage was limited. Two Semantic Scholar queries failed on rate limits, and the unavailable arXiv API was replaced by the site search. OpenAlex relevance ranking, which weights citations, favored highly cited work from other fields. Several of the closest works surfaced only through citation, and theses, non-English literature and survey methodology were not searched.
% R58
Within this search and the works read, we did not identify a study that estimates the minimum partner-drift strength at which individual panel trajectories become distinguishable from a family of matched no-network persistence nulls as a function of the response channel's signal share; absence from a limited search is not proof of absence.
% R59
No single component is new: minimum signal against persistence \cite{moor2026,coding2024}, reliability-dependent sensitivity \cite{mulder2023,dang2026}, power for influence \cite{stadtfeld2020} and detectability maps \cite{urena2026} all have antecedents. The proposed extension combines individual trajectories on an unobserved network, a separate response channel with known $r$, persistence nulls matched on short-lag moments, and a fixed panel-level discriminator with a censored strength boundary.
% R74
The conjunction serves a specific comparison. Moor-Smith and Carpentras ask when drift improves prediction of repeated population distributions over no change. Soni and colleagues compare event predictions on an observed network; longitudinal measurement studies separate response instability from persistent change. Here the question is what longer-lag information remains in individual panels after persistence rivals reproduce short-term stability, and how that information changes under a separately controlled observation channel. Short-lag matching and fixed marginal distributions direct the comparison toward longer-lag differences in trajectories. The boundary then describes the reach of one declared discriminator under these conditions. This is a bounded extension for evaluating rival accounts of attitude change; it does not establish priority, an optimal detection limit or causal influence.
